% !TeX program = xelatex
\documentclass[12pt,a4paper]{article}
\usepackage[margin=2cm]{geometry}
\usepackage{fontspec}
\usepackage{amsmath,amssymb}
\usepackage{enumitem}
\setmainfont{Times New Roman}
\setlength{\parindent}{0pt}
\setlength{\parskip}{6pt}

\begin{document}
\begin{center}
  {\Large\bfseries ĐÁP ÁN ĐỀ THI TUYỂN SINH VÀO LỚP 10 THPT}\\
  {\bfseries Thành phố Hà Nội, năm học 2008--2009}\\
  {\itshape Phần Đại số}
\end{center}

\textbf{Câu 1.}

Điều kiện xác định: $x>0$.

\begin{enumerate}[label=\alph*)]
  \item Ta có
  \begin{align*}
  P&=\left(\frac{1}{\sqrt{x}}+\frac{\sqrt{x}}{\sqrt{x}+1}\right)
  \cdot\frac{x+\sqrt{x}}{\sqrt{x}}\\
  &=\frac{x+\sqrt{x}+1}{\sqrt{x}(\sqrt{x}+1)}
  \cdot\frac{\sqrt{x}(\sqrt{x}+1)}{\sqrt{x}}\\
  &=\frac{x+\sqrt{x}+1}{\sqrt{x}}.
  \end{align*}

  \item Với $x=4$,
  \[
  P=\frac{4+2+1}{2}=\frac{7}{2}.
  \]

  \item Ta có
  \begin{align*}
  \frac{x+\sqrt{x}+1}{\sqrt{x}}&=\frac{13}{3}\\
  3x-10\sqrt{x}+3&=0\\
  (\sqrt{x}-3)(3\sqrt{x}-1)&=0.
  \end{align*}
  Do đó $x=9$ hoặc $x=\dfrac{1}{9}$; cả hai giá trị đều thỏa mãn điều kiện.
\end{enumerate}

\textbf{Câu 2.}

Gọi $x,y$ lần lượt là số chi tiết máy tổ I và tổ II sản xuất trong tháng thứ nhất ($x,y>0$). Khi đó
\[
\begin{cases}
x+y=900,\\
1{,}15x+1{,}10y=1010.
\end{cases}
\]
Tương đương
\[
\begin{cases}
x+y=900,\\
23x+22y=20200.
\end{cases}
\]
Giải hệ được $x=400$, $y=500$.

Vậy trong tháng thứ nhất, tổ I sản xuất 400 chi tiết máy và tổ II sản xuất 500 chi tiết máy.

\textbf{Câu 3.}

Hoành độ giao điểm của $(P)$ và $(d)$ thỏa mãn
\[
\frac{x^2}{4}=mx+1
\quad\Longleftrightarrow\quad
x^2-4mx-4=0. \tag{1}
\]

\begin{enumerate}[label=\alph*)]
  \item Phương trình (1) có
  \[
  \Delta'=4m^2+4>0
  \]
  với mọi $m$. Vì vậy $(d)$ luôn cắt $(P)$ tại hai điểm phân biệt $A,B$.

  \item Gọi $x_1,x_2$ là hai nghiệm của (1). Khi đó
  \[
  x_1+x_2=4m,\qquad x_1x_2=-4.
  \]
  Với $A\left(x_1,\dfrac{x_1^2}{4}\right)$ và
  $B\left(x_2,\dfrac{x_2^2}{4}\right)$, ta có
  \begin{align*}
  S_{AOB}
  &=\frac12\left|x_1\frac{x_2^2}{4}-x_2\frac{x_1^2}{4}\right|\\
  &=\frac18\left|x_1x_2(x_2-x_1)\right|.
  \end{align*}
  Mặt khác,
  \[
  |x_2-x_1|
  =\sqrt{(x_1+x_2)^2-4x_1x_2}
  =4\sqrt{m^2+1}.
  \]
  Suy ra
  \[
  \boxed{S_{AOB}=2\sqrt{m^2+1}}.
  \]
\end{enumerate}

\textbf{Câu 5.}

Đặt $a=x-2$. Khi đó $x-1=a+1$ và $x-3=a-1$. Ta có
\begin{align*}
Q&=(a+1)^4+(a-1)^4+6(a+1)^2(a-1)^2\\
 &=8a^4+8\ge 8.
\end{align*}
Dấu bằng xảy ra khi $a=0$, tức là $x=2$.

Vậy giá trị nhỏ nhất của $Q$ là $\boxed{8}$.

\end{document}
